\section{Discussion}
\label{sec:discussion}

We have described a certified triage instrument for combinatorial perturbation screens. It treats
a screen's single-gene effect vectors as a dictionary and decides, by convex geometry, whether a
measured combination reaches a cell-state direction that the singles cannot---answering ``is this
combination emergent?'' with a Farkas/KKT certificate rather than a prediction, and answering
``which unmeasured combinations are worth running?'' with a training-free score computed from the
singles alone. The value of the framework is not a new list of synergistic gene pairs---those are
already in the source screens and we use them only as positive controls---but a new \emph{kind of
output}: a falsifiable, noise-aware emergence certificate computed directly from measured effects,
which no forward predictor supplies.

\subsection{Interpretation}
The central methodological finding is that predictive accuracy is the wrong axis for claiming a
combination is special. A self-contained learned model and a one-line additive rule predict
held-out double effects equally well ($0.896$ vs.\ $0.897$), consistent with a wave of
2025--2026 benchmarks \citep{AhlmannEltze2025,Wong2025,Kernfeld2025,Atti2025}; a screening team
that ranked combinations by predicted accuracy would be ranking by a quantity on which the best
available model ties a trivial baseline. The convex-cone view instead asks a decision question
with a checkable answer on both sides, and only it returns a certificate: on all $131$ Norman
doubles the measured effect leaves the single-gene cone with a separating hyperplane at machine
precision.

That geometric answer becomes trustworthy only once it clears a noise bar. The raw unreachable
fraction is confounded by effect magnitude (Spearman $-0.56$), because a normalized residual
inflates for low-magnitude combinations measured with the same absolute noise; the
noise-injection null removes the confound ($+0.14$) and yields a two-bar verdict that separates
``emergent at all above noise'' ($129/131$) from ``emergent by a large margin relative to its own
noise'' ($35/131$). We report both bars because they answer different questions, and because a
learned point predictor, emitting no null, can supply neither. Reassuringly, the certificate is
not a measurement artefact: it recovers classical non-additivity at partial Spearman $0.62$ after
controlling for magnitude, so the pairs it certifies are the biologically synergistic ones.

The complement to the certificate is prospective, and here a single scalar carries the signal.
The negative cosine between two single effects ranks emergent combinations $2.4\times$ above base
rate without any measurement of the combination---dissimilar singles are more likely to combine
into something the cone cannot reach. This is a training-free rule a team can apply to a full
combinatorial catalogue before committing a single well.

\subsection{Limitations}
\label{sec:limitations}
Four limitations bound the conclusion, and each maps to a concrete next step.

The most consequential is that emergence is defined \emph{relative to the measured dictionary}
under an additive model, not as biological synergy in an absolute sense. A combination we certify
as emergent reaches a state no non-negative mixture of the \emph{measured} singles reaches; if the
true single-gene effect is mis-estimated, or a relevant single is absent from the library, the
verdict changes. The $k$-way analysis makes this dependence explicit and turns it into a feature:
enlarging the reference set from singles to singles-plus-doubles shrinks the certified set from
$40$ to $16$, because emergence is a property of the reference, and the method forces that
reference to be stated.

Second, the higher-order validation is uneven. The order-$2$ certificate rests on real measured
doubles; the order-$3$ certificate is validated only on a synthetic planted-epistasis screen,
because the Norman screen measures no triples. The AUC $1.00$ on synthetic triples establishes
that the code path discriminates planted $3$-way emergence from additive and $2$-way-reducible
combinations---it is a validation of the machinery, not evidence about real biological $3$-way
epistasis, which no available data can supply.

Third, generalization is partial and we state the boundary. The certificate and its noise-aware
discipline transfer to an orthogonal screen (CaRPool-seq; Cas13d knockdown in THP-1): the
magnitude confound and its noise-aware correction reproduce, and the two-bar certificate holds
(CaRPool $48\%$ certified). But the prospective triage rule transfers only in rank against the raw
label ($\sim\!2.2\times$ in both screens); its magnitude-controlled component survives in Norman
and is not significant in CaRPool, so the triage score should be recalibrated per screen rather
than assumed portable.

Finally, external validity is bounded by provenance and by the transcriptomic proxy. The Norman
combinatorial file used here is labelled A549 whereas the canonical Norman 2019 line is K562, a
discrepancy we carry through every result; and matching a target in expression space is not the
same as achieving the corresponding cellular phenotype---functional confirmation remains a
separate wet-lab endpoint.

\subsection{Outlook}
The most valuable next experiments follow directly from the limitations. Measuring a modest panel
of triple perturbations would replace the synthetic order-$3$ validation with an in-domain one and
test whether the certified-set shrinkage seen from singles to doubles continues to higher order.
Applying the training-free triage prospectively---selecting the top-ranked unmeasured combinations
in a fresh screen and measuring the realized emergence rate against random selection---would
convert the retrospective $2.4\times$ enrichment into a forward-validated hit rate. And extending
the dictionary across more screens and modalities would turn the two-screen transfer result into a
generalization curve, and calibrate when the prospective score is portable. In each case the
framework's output is intended to make combinatorial screening more efficient by exposing, before
the wells are run, which combinations can reach states the singles cannot.

\paragraph{A causal-inference agenda.}
Read as a causal-inference method, the certificate is a design-based counterfactual query: the
single-gene effect vectors are average treatment effects from a randomized perturbation
experiment, and an emergence verdict asks whether a combination's measured effect exceeds what any
non-negative mixture of those single treatment effects produces. That framing makes the identifying
assumptions explicit and turns each into a concrete robustness check. Four are worth naming. First,
\emph{interference} (SUTVA): pooled Perturb-seq lets a perturbed cell's secreted signals alter a
neighbour's state, which would bias the single-gene effects the cone is built from; an
arrayed-versus-pooled or MOI-titration comparison would bound this spillover. Second,
\emph{measurement noise as coordinated bias}: our noise-injection null models undirected per-gene
noise, but a coordinated attenuation of the combination-specific signal would be a different failure
mode---a sensitivity radius on the certified verdict quantifies how much coordinated bias would
overturn it. Third, \emph{construct validity}: negative-control combinations (pairs expected to be
additive) should not certify, and the specificity of the synthetic screen ($0/80$ false positives)
is the in-silico version of that check. Fourth, \emph{guide non-compliance}: incomplete knockdown
or activation attenuates the measured effect, an errors-in-variables problem an instrumental-variables
treatment of guide efficiency would address. The computable half of this agenda---sensitivity radii,
negative-control outcomes, and signed construct validity for the emergence certificate---is released
as a causal-validation dossier accompanying the code; the wet-lab experiments above are its
continuation.
